\documentclass{crypto-exercise}
\usepackage{amsthm}
\author{Sven Laur}
\editor{Sven Laur}
\tags{computational indistinguishability, distributions, constructive hybrid argument}

\renewcommand{\ADVIND}[2]{\ADV^{\mathsf{ind}}_{#1}(#2)}
\newcommand{\ts}{t_{\mathrm{s}}}

\begin{document}
\begin{exercise}{Constructive hybrid argument}
Let $\XXX_0$ and $\XXX_1$ efficiently samplable distributions that are $(t, \varepsilon)$-indistinguishable. Show that distributions $\XXX_0$ and $\XXX_1$ remain computationally indistinguishable even if the adversary can get $n$ samples.
As the first step, estimate computational distances between following games
\begin{align*}
  & \begin{game}{\GAME_{00}^\AD}
    &  x_0 \gets \XXX_0 \\
    &  x_1 \gets \XXX_0 \\
    & \RETURN \AD(x_0,x_1)
  \end{game} 
&
 & \begin{game}{\GAME_{01}^\AD}
    &  x_0 \gets \XXX_0 \\
    &  x_1 \gets \XXX_1 \\
    & \RETURN \AD(x_0,x_1)
  \end{game} 
&
 & \begin{game}{\GAME_{11}^\AD}
    &  x_0 \gets \XXX_1 \\
    &  x_1 \gets \XXX_1 \\
    & \RETURN \AD(x_0,x_1)
  \end{game} 
&
\end{align*}
and then generalise the argumentation to the case, where the adversary $\AD$ gets $n$ samples from a distribution $\XXX_i$.  Prove the statement by providing a \textbf{single} adversary $\ADB^\AD$ that plays against  indistinguishability games
\begin{align*}
  & \begin{game}{\BGAME_{0}^\ADB}
    &  x \gets \XXX_0 \\
    & \RETURN \ADB(x)
  \end{game} 
&
 & \begin{game}{\BGAME_{1}^\ADB}
    &  x \gets \XXX_1 \\
    & \RETURN \ADB(x)
  \end{game} 
\end{align*}
Estimate how the advantage of $\advIndXX{\GAME_0,\GAME_1}{\AD}$ and $\advIndXX{\BGAME_0,\BGAME_1}{\ADB}$ are related.
Generalise the construction to general case. Why do we need to assume that distributions $\XXX_0$ and $\XXX_1$ are efficiently samplable?
\end{exercise}


\begin{solution}
\end{solution}
\end{document}